\begin{frame}{Aufgabe}
	% WiSe 10/11 Aufgabe 6 c
	Die Sprache $L\subseteq \{a,b\}^* $ sei definiert als die Menge aller Wörter $w$, die folgende Bedingungen erfüllen 
	\begin{align*}
		N_b(w) &> N_a(w)  \\ \forall v_1,v_2 \in \{a,b\}^* : w &\neq v_1 bb v_2 
	\end{align*}

	Geben Sie einen endlichen Akzeptor an, der $L$ erkennt. \\
	
	\bigskip
	\pause
	Hinweis: Umgangssprachliche Formulierung der Sprache:\\ \pause
	Anzahl b \textbf{echt} größer als Anzahl a \textbf{und}\\
	Nirgends zwei b hintereinander\\
	
	\medskip
	\pause
	Hinweis 2: Was passiert, wenn $w$ mit einem a startet?\\ \pause
	$w$ kann nicht mehr in $L$ liegen (wegen der 2. Bedingung)
\end{frame}

\begin{frame}{Lösung}
	\begin{figure}
		\centering
		\includegraphics[width=0.6\linewidth]{automaten/Loesung2.pdf}
	\end{figure}
\end{frame}

% TODO: Evtl. nächstes Tut
\begin{frame}
	Gibt es einen endlichen Akzeptor für die Sprache $$L = \{ a^n b^n | n \in \N_0 \}$$
	\pause
	Nein! Siehe VL!
	% TODO: SOON Hinweis auf TM zum erkennen
\end{frame}